\chapter{Expériences en conditions réelles}\label{ch:rs:live-experiments}

Un nouvel algorithme d'exécution traite dix pour cent des ordres de la société pendant deux semaines, soit 2\,485 sur 24\,000, et économise
en moyenne 0.55 point de base d'\omterm{def:rs:simulation-versus-live:calibration}{implementation shortfall}, avec une erreur type de 0.48. L'intervalle à 95~\% va de $-1.49$ à 0.39~: l'expérience ne permet pas de distinguer une économie d'un coût. Avec la moitié du flux, il
faudrait 77 jours de trading pour détecter une économie de 0.3 point de base, et 38 avec la réduction de variance de ce chapitre. Même une réponse parfaite porterait sur une
question plus étroite que celle du desk~: dans le flux simulé du chapitre, le nouvel algorithme économise 0.30 point de base sur ses propres ordres, mais ajoute 0.25 point de base aux autres ordres de
la société sur la même valeur, du même côté et le même jour, si bien que basculer tout le flux économise 0.08. La recherche atteint la production par le \omterm{def:rs:live-experiments:paper}{trading simulé}, les \omterm{def:rs:live-experiments:canary}{déploiements canari}
et progressifs, puis elle est mesurée sur le flux réel~; ce chapitre traite des deux, avec \texttt{firm.abtest}.

\section{Trading simulé et trading fantôme}

\begin{definition}[Trading simulé]\label{def:rs:live-experiments:paper}
Le \emph{trading simulé}\index{trading simulé} (paper trading) fait tourner une nouvelle stratégie en temps réel sur des données de marché en direct~: il calcule ses signaux et génère ses ordres, qui sont
enregistrés et contrôlés, mais jamais envoyés, et les exécute au moyen d'un \omterm{def:rs:event-driven-backtests:fill}{modèle d'exécution} (fill model).
\end{definition}

Le \omterm{def:rs:live-experiments:paper}{trading simulé} teste tout ce qui ne dépend pas de la réaction du marché. Les flux de données en direct arrivent avec leur calendrier réel et leurs lacunes~; les signaux calculés
en direct peuvent être comparés aux mêmes signaux recalculés par le code de recherche sur les mêmes données, valeur par valeur~; les ordres réussissent ou échouent aux contrôles de risque
pré-trade~; la stratégie traverse le calendrier opérationnel (démarrage et arrêt, opérations sur titres, jours fériés, rééquilibrage de l'indice)
qu'un \omterm{def:rs:vectorised-backtests:backtest}{backtest} réduit à des lignes. Ce que le \omterm{def:rs:live-experiments:paper}{trading simulé} ne peut pas tester, ce sont les exécutions~: c'est un simulateur tournant sur des données en direct, avec les angles morts de tout simulateur~: le chapitre~19
a montré qu'un rejeu sans les ordres propres de la stratégie manque les transactions qui s'arrêtaient sur eux et les niveaux où elles étaient seules. Le \omterm{def:rs:simulation-versus-live:reconciliation}{trading fantôme}
(chapitre~19) compare le simulateur au trading en direct d'une stratégie existante~; le \omterm{def:rs:live-experiments:paper}{trading simulé} est l'étape qui précède le premier trading en direct d'une stratégie nouvelle.

La condition de sortie du \omterm{def:rs:live-experiments:paper}{trading simulé} n'est donc pas son P\&L. C'est la parité~: les signaux en direct égaux aux signaux de recherche à une tolérance énoncée près, chaque
ordre qu'il aurait envoyé soit accepté par les contrôles de risque, et qu'un cycle opérationnel complet soit franchi sans correction manuelle. Une stratégie au P\&L simulé satisfaisant, dont
les signaux en direct diffèrent de ceux de la recherche, n'a pas été testée~; elle a été remplacée par une autre.

\section{Déploiements canari et déploiements progressifs}

\begin{definition}[Déploiement canari, déploiement progressif]\label{def:rs:live-experiments:canary}
Un \emph{déploiement canari}\index{déploiement canari} fait tourner une nouvelle version d'un composant de trading sur une petite part du flux de production (le canari) en parallèle
avec l'ancienne version, et compare les deux avant que la nouvelle n'en reçoive davantage. Un \emph{déploiement progressif}\index{déploiement progressif} augmente la part par paliers (par exemple
1~\%, 10~\%, 50~\%, 100~\%), chaque palier étant conditionné par des contrôles énoncés, avec un chemin de retour testé vers l'ancienne version à chaque palier.
\end{definition}

La défaillance de déploiement la plus coûteuse du dossier public était un \omterm{def:rs:live-experiments:canary}{déploiement progressif}. À partir du 27 juillet 2012, Knight Capital a placé un nouveau code de routage d'ordres sur
un nombre limité de serveurs, sur des jours successifs, avant un nouveau programme de la NYSE qui démarrait le 1 août. Un technicien n'a pas copié le nouveau code sur
un des huit serveurs~; aucune seconde personne n'a relu le déploiement, et aucune procédure écrite ne l'exigeait. Le 1 août, le huitième serveur exécutait une ancienne fonction défectueuse que
le drapeau réaffecté du nouveau code réactiva~: pour 212 ordres parents entrants, il envoya des millions d'ordres enfants, 4 millions d'exécutions sur 154 valeurs, pour plus de
397 millions d'actions en environ 45 minutes. La société s'est retrouvée avec une position acheteuse nette d'environ 3.5~milliards de dollars sur 80 valeurs et une position vendeuse nette d'environ 3.15~milliards de dollars sur
74, et a enregistré une perte de 460~millions de dollars. Avant l'ouverture, ses systèmes avaient envoyé 97 courriels automatiques citant l'ancienne fonction par son nom~; ils n'avaient pas été conçus comme
des alertes, et personne ne les lisait. L'ordonnance de la SEC d'octobre 2013 a imposé une amende de 12~millions de dollars.

Trois leçons en découlent, aucune de nature statistique. Une étape ne vaut que par la vérification qu'elle fait de l'état déployé~: chaque serveur, chaque processus et chaque configuration exécute la
version que l'étape annonce. Une alerte est un message conçu et acheminé pour donner lieu à une action, et non une ligne de journal qui se trouve être envoyée par courriel. Une étape doit être plafonnée
dans ce qu'elle peut faire~: un canari portant 1~\% du flux doit détenir 1~\% des limites de risque, et les invariants tels que «~la quantité enfant ne dépasse jamais celle du parent~» doivent être
vérifiés sur chaque ordre, où ils se déclenchent dès la première violation. La comparaison statistique du canari et du témoin relève d'une autre classe de défaut.

\begin{proposition}[Durée d'un canari statistique]\label{prop:rs:live-experiments:canary}
Un canari reçoit une part $s$ d'un flux~; un défaut ajoute $\delta$ à une mesure par ordre d'écart type $\sigma$~; l'alarme est un test $z$ unilatéral de la
moyenne du canari par rapport à celle du témoin, au seuil $z$. L'alarme attendue survient après environ
\[
n \approx \frac{z^2\sigma^2}{\delta^2(1 - s)}
\]
ordres du canari, quantité presque indépendante de la part lorsqu'elle est faible.
\end{proposition}

\begin{proof}
Après $n$ ordres du canari, le témoin en compte $m = n(1-s)/s$, et la différence des moyennes a pour variance $\sigma^2/n + \sigma^2/m = \sigma^2/(n(1-s))$. La statistique espérée
$\delta\sqrt{n(1-s)}/\sigma$ atteint $z$ pour le $n$ indiqué.
\end{proof}

Avec le flux du chapitre ($\sigma = 23.1$ points de base, 2\,400 ordres par jour), un défaut coûtant 5 points de base par ordre et une alarme à $z = 3$, un canari à 1~\%
déclenche l'alarme après 194 ordres touchés et 8.1 jours, un canari à 10~\% après 214 ordres et 0.89 jour, un canari à 50~\% après 385 ordres et 0.32 jour. Un canari plus petit
ne réduit pas les dégâts subis avant une alarme statistique~; il les étale sur davantage de temps. Ce qu'un petit canari achète, c'est une faible exposition aux défauts
détectés par les invariants et les limites, et du temps pour que des personnes examinent la situation.

\begin{omfigure}
\begin{tikzpicture}[stage/.style={draw=omDef,fill=omDef!8,rounded corners=2pt,text width=1.85cm,align=center,font=\scriptsize,minimum height=1.1cm},
  gate/.style={text width=1.95cm,align=center,font=\scriptsize,text=black!70},arr/.style={-{Stealth},thick,omDef}]
\node[stage] (r) at (0,0) {recherche\\backtests};
\node[stage] (p) at (2.35,0) {trading simulé\\(sans ordres)};
\node[stage] (c) at (4.7,0) {canari\\1~\% du flux};
\node[stage] (a) at (7.05,0) {test A/B\\par valeur-jour};
\node[stage] (h) at (9.4,0) {50~\%};
\node[stage] (f) at (11.75,0) {déploiement complet};
\foreach \x/\y in {r/p,p/c,c/a,a/h,h/f} \draw[arr] (\x) -- (\y);
\node[gate,below=3pt of r] {PBO, SR déflaté, modèle de coûts (ch.~20, 27)};
\node[gate,below=3pt of p] {parité des signaux, contrôles de risque, calendrier complet};
\node[gate,below=3pt of c] {état déployé vérifié, invariants, 1~\% des limites};
\node[gate,below=3pt of a] {effet et puissance, indicateurs de garde-fou};
\node[gate,below=3pt of h] {surveillance, capacité (ch.~28)};
\node[gate,below=3pt of f] {kill switch, chemin de retour testé};
\draw[omThm,thick,-{Stealth},dashed] (f.north) to[bend right=12] node[above,font=\scriptsize,text=omThm] {retour arrière à toute étape} (c.north);
\end{tikzpicture}

\omcaption{De la recherche à la production~: chaque étape est conditionnée par des contrôles qu'elle peut réellement effectuer. Le \omterm{def:rs:live-experiments:paper}{trading simulé} ne mesure pas les exécutions~; les statistiques d'un canari
ne permettent pas de rattraper à temps un ordre hors de contrôle~; ses invariants et ses limites, eux, le peuvent.}\label{fig:rs:live-experiments:stages}
\end{omfigure}

\section{Tests A/B randomisés sur le flux réel}

\begin{definition}[test A/B, unité de randomisation, interférence, déséquilibre du ratio d'échantillonnage]\label{def:rs:live-experiments:ab}
Un \emph{test A/B}\index{test A/B} compare deux versions d'un composant en affectant au hasard des unités du flux réel à chacune et en comparant un résultat.
L'\emph{unité de randomisation}\index{unité de randomisation} désigne ce qui est affecté~: un ordre, un couple valeur-jour, un jour, un client. \emph{Interférence}\index{interférence}
survient lorsque le résultat d'une unité dépend de l'affectation des autres, si bien que la différence entre les bras d'un test diffère de l'effet de
basculer chaque unité. Un \emph{déséquilibre du ratio d'échantillonnage}\index{déséquilibre du ratio d'échantillonnage} (sample ratio mismatch, SRM) est une répartition observée entre les bras, qui s'écarte,
au-delà du hasard, de celle qui était prévue (Fabijan et coauteurs).
\end{definition}

Le flux simulé du chapitre compte 2\,400 ordres parents par jour sur 500 valeurs, tirées avec une popularité inversement proportionnelle au rang. Pour chaque ordre,
l'\omterm{def:rs:simulation-versus-live:calibration}{implementation shortfall}, exprimée en points de base par rapport à son \omterm{def:rs:simulation-versus-live:calibration}{prix d'arrivée}, est l'estimation de coût pré-trade de la société (écart type 6.9 entre ordres), l'erreur de cette estimation,
le mouvement de marché ajusté du bêta sur la fenêtre planifiée de l'ordre, signé selon son sens, du bruit idiosyncratique, et un effet jour commun à tous les ordres~; la moyenne
vaut 11.7 et l'écart type 23.1. Le nouvel algorithme économise 0.30 point de base sur chaque ordre qu'il traite, et ajoute 0.25 aux autres ordres de la même valeur,
même côté et même jour lorsque tous sont sur le nouvel algorithme (il exécute plus tôt dans la fenêtre et laisse à ses voisins davantage d'impact).

\paragraph{Affectation.} \texttt{firm.abtest.assign} hache l'identifiant de l'unité avec le nom de l'expérience et compare les huit premiers octets de l'empreinte SHA-256
empreinte, lue comme une fraction, à la part prévue. L'affectation est déterministe (un même ordre va toujours dans le même bras, sur chaque serveur, sans une
table partagée), indépendante entre expériences de noms différents, et reproductible a posteriori par quiconque détient les identifiants. Dans l'expérience de deux semaines
d'introduction, 2\,485 ordres sur 24\,000 sont allés au nouvel algorithme contre 2\,400 prévus~: $z = 1.83$, une p-value de 0.07, aucun déséquilibre à 5~\%. Un
déséquilibre ne serait pas une simple curiosité statistique~: c'est le symptôme d'ordres perdus, dupliqués ou réacheminés dans un bras, et il invalide la comparaison tant
qu'il n'est pas expliqué.

\paragraph{\omterm{def:rs:live-experiments:ab}{Interférence}.} Dans le flux, 88~\% des ordres partagent leur action, leur sens et leur jour avec un autre ordre de la société. Lorsque les ordres sont randomisés un à
un, un ordre traité et un ordre de contrôle ont, en moyenne, la même part de frères traités, de sorte que la différence entre les bras mesure
l'économie directe, $-0.30$. La question du desk est de savoir ce qui se passe lorsque chaque ordre bascule~: $-0.30 + 0.25 \times 0.883 = -0.08$. Une \omterm{def:rs:live-experiments:ab}{unité de randomisation} doit contenir
l'\omterm{def:rs:live-experiments:ab}{interférence}~: ici le couple action-jour, c'est-à-dire tous les ordres d'une action un jour donné dans le même bras. Le jour lui-même contient toutes les \omterm{def:rs:live-experiments:ab}{interférences} d'une journée, mais
n'offre qu'une seule unité par jour, et l'effet du jour devient alors le bruit~; alterner le traitement dans le temps de cette façon constitue un plan en bascule (Bojinov, Simchi-Levi
et Zhao), courant lorsque les unités interfèrent, au prix d'un petit nombre d'unités et d'un report d'une période à la suivante.

\omcode{code/firm/abtest/firm_abtest.py}{74}{91}{La comparaison lorsque des grappes entières sont affectées~: moyennes pondérées par les ordres, erreur type robuste aux grappes.}\label{lst:rs:live-experiments:cluster}

\section{Réduction de variance}

\begin{definition}[CUPED]\label{def:rs:live-experiments:cuped}
La méthode \emph{CUPED}\index{CUPED} (expériences contrôlées utilisant des données antérieures à l'expérience~; Deng, Xu, Kohavi et Walker) compare les bras sur le résultat ajusté
$y - (x - \bar x)^\top\theta$, où $x$ sont des covariables que le traitement ne peut pas affecter et $\theta$ est le coefficient de la régression de $y$ sur $x$ regroupée sur
les deux bras.
\end{definition}

\begin{proposition}[Ce que CUPED apporte]\label{prop:rs:live-experiments:cuped}
La différence ajustée des moyennes est sans biais pour l'effet du traitement, et sa variance vaut la variance non ajustée multipliée par $1 - R^2$, où $R^2$ est la part de
la variance du résultat expliquée par les covariables.
\end{proposition}

\begin{proof}
La randomisation donne aux covariables la même loi dans les deux bras, donc $\mathbb{E}[\bar x_T - \bar x_C] = 0$ et l'ajustement ne retire rien à la
différence espérée. La variance de $y - x^\top\theta$ est minimale au coefficient de régression, où elle vaut $\operatorname{Var}(y)(1 - R^2)$~; la
différence de deux moyennes indépendantes hérite du facteur. Estimer $\theta$ sur les données regroupées modifie cela d'un terme d'ordre $1/n$.
\end{proof}

Les covariables doivent être fixées avant l'affectation ou se situer hors de portée de la société. L'estimation de coût pré-trade convient, et explique 9~\% de la variance de
l'\omterm{def:rs:simulation-versus-live:calibration}{implementation shortfall}. Il en va de même pour le mouvement de marché ajusté du bêta sur la fenêtre planifiée de l'ordre~: le gérant de portefeuille fixe la fenêtre avant
l'affectation de l'ordre, et le marché ne bouge pas parce qu'un des algorithmes de la société plutôt que l'autre le traite. Avec les deux covariables, $R^2$ vaut 0.51, proche
de la division par deux rapportée par Deng et ses coauteurs pour les expériences de Bing. La durée réalisée de l'ordre, son taux d'exécution ou sa participation ne conviennent pas~:
l'algorithme les modifie, et les ajuster retire une partie de l'effet. La post-stratification (Miratrix, Sekhon et Yu) est \omterm{def:rs:live-experiments:cuped}{CUPED} avec des indicatrices de strates comme
covariables, presque aussi efficace qu'une affectation par strates effectuée à l'avance~; pour un plan randomisé par couple action-jour, l'analyse au sein de chaque jour (en retranchant à chaque jour
sa moyenne du résultat et des covariables) est la stratification qui supprime l'effet du jour.

Chaque plan présenté ci-dessous a été exécuté sur 100 expériences simulées de 60 jours sur la moitié du flux (l'estimation moyenne, la dispersion des estimations, la moyenne
des erreurs types rapportées, et les jours qu'implique cette erreur type pour une puissance de 80~\% sur 0.3 point de base~:

\begin{center}
\small
\begin{tabular}{@{}lrrrr@{}}
\toprule
plan et estimateur & estimation & dispersion & erreur type & jours pour 0.3 pb\\
\midrule
par ordre~: différence des moyennes & $-0.30$ & 0.123 & 0.122 & 77\\
par ordre~: stratifié (quintiles pré-trade) & $-0.30$ & 0.119 & 0.117 & 72\\
par ordre~: \omterm{def:rs:live-experiments:cuped}{CUPED}, estimation pré-trade & $-0.30$ & 0.121 & 0.116 & 70\\
par ordre~: \omterm{def:rs:live-experiments:cuped}{CUPED}, pré-trade et marché & $-0.29$ & 0.095 & 0.085 & 38\\
par jour & $-0.16$ & 0.656 & 0.788 & 3\,253\\
par couple action-jour & $-0.09$ & 0.235 & 0.235 & 290\\
par couple action-jour, au sein du jour & $-0.09$ & 0.188 & 0.184 & 178\\
par couple action-jour, au sein du jour, \omterm{def:rs:live-experiments:cuped}{CUPED} & $-0.08$ & 0.093 & 0.082 & 35\\
\bottomrule
\end{tabular}
\end{center}

Tout plan au niveau de l'ordre estime l'économie directe, $-0.30$~; tout plan qui affecte des couples action-jour ou des jours entiers estime celle du déploiement, $-0.08$ (le plan par jour,
$-0.16$, se situe dans sa dispersion de 0.66). Les erreurs types rapportées correspondent à la dispersion des estimations d'une expérience à l'autre. \omterm{def:rs:live-experiments:cuped}{CUPED} avec la covariable de marché
divise par deux le nombre de jours nécessaires~; l'estimation pré-trade seule et la stratification sur celle-ci apportent peu, car celle-ci explique peu. Randomiser par
jour multiplie par 42 le nombre de jours. Randomiser par couple action-jour et analyser au sein du jour, avec \omterm{def:rs:live-experiments:cuped}{CUPED}, est aussi précis que le meilleur plan au niveau de l'ordre et
répond à la bonne question~; la réponse est petite, et détecter une économie de 0.08 point de base avec une puissance de 80~\% prendrait 503 jours de négociation.

\section{Dimensionnement et arrêt d'une expérience}

\begin{definition}[Effet minimal détectable, test séquentiel]\label{def:rs:live-experiments:mde}
L'\emph{effet minimal détectable}\index{effet minimal détectable} d'un plan est le plus petit effet réel qu'il détecte avec une puissance annoncée à une taille annoncée.
Un \emph{test séquentiel}\index{test séquentiel} examine les données à mesure qu'elles arrivent et peut s'arrêter à tout moment, sa taille étant contrôlée sur l'ensemble des examens que
il pourrait effectuer.
\end{definition}

\begin{proposition}[Effet minimal détectable et taille d'échantillon]\label{prop:rs:live-experiments:mde}
Pour un test bilatéral de taille $\alpha$ et de puissance $1-\beta$ sur une estimation d'erreur type $\mathrm{se}$, l'\omterm{def:rs:live-experiments:mde}{effet minimal détectable} vaut
$(z_{1-\alpha/2} + z_{1-\beta})\,\mathrm{se}$, soit $2.80\,\mathrm{se}$ à 5~\% et 80~\%. Avec un écart type du résultat $\sigma$ et une part $s$ de $n$
unités traitées, détecter un effet $\Delta$ exige
\[
n = \Bigl(\frac{(z_{1-\alpha/2} + z_{1-\beta})\,\sigma}{\Delta}\Bigr)^2\Bigl(\frac1s + \frac1{1-s}\Bigr).
\]
\end{proposition}

\begin{proof}
Le test rejette lorsque $|\hat\Delta|/\mathrm{se} > z_{1-\alpha/2}$~; si $\hat\Delta \sim N(\Delta, \mathrm{se}^2)$, cela a pour probabilité $1-\beta$ (en négligeant
l'autre queue) lorsque $\Delta/\mathrm{se} = z_{1-\alpha/2} + z_{1-\beta}$. On remplace $\mathrm{se}^2 = \sigma^2/(sn) + \sigma^2/((1-s)n)$.
\end{proof}

L'écart type de 23.1 points de base et 2\,400 ordres par jour donnent 77.7 jours à la moitié du flux pour 0.3 point de base, et 38.2 avec l'écart type résiduel
de \omterm{def:rs:live-experiments:cuped}{CUPED} de 16.2~; les expériences simulées, qui rejettent dans 78~\% des cas à 80 jours sans \omterm{def:rs:live-experiments:cuped}{CUPED} et 76~\% à 40 jours avec, confirment la formule
(la \cref{fig:rs:live-experiments:power}). La répartition égale minimise la taille de l'échantillon~; une expérience à 10~\%, comme celle de l'introduction, demande $(1/0.1 + 1/0.9)/4 = 2.8$
fois plus d'ordres.

\begin{omfigure}
\begin{tikzpicture}
\begin{axis}[width=11cm,height=5.6cm,xlabel={jours de négociation à la moitié du flux},ylabel={puissance pour 0.3 pb},tick label style={font=\small},label style={font=\small},
  xmin=0,xmax=120,ymin=0,ymax=1.02,grid=major,grid style={gray!25},legend columns=3,
  legend style={font=\footnotesize,at={(0.5,-0.3)},anchor=north,draw=none,fill=none,/tikz/every even column/.append style={column sep=8pt}},
  table/col sep=comma]
\addplot[black!60,thick] table[x=days,y=raw]{figdata/research/21-live-experiments/power.csv};
\addplot[omDef,very thick] table[x=days,y=cuped]{figdata/research/21-live-experiments/power.csv};
\addplot[omThm,thick,dashed] table[x=days,y=day]{figdata/research/21-live-experiments/power.csv};
\addplot[only marks,mark=o,mark size=2pt,black!60,forget plot] table[x=days,y=raw]{figdata/research/21-live-experiments/power_sim.csv};
\addplot[only marks,mark=*,mark size=1.8pt,omDef,forget plot] table[x=days,y=cuped]{figdata/research/21-live-experiments/power_sim.csv};
\draw[black!50,densely dotted] (axis cs:0,0.8) -- (axis cs:120,0.8);
\legend{par ordre,par ordre avec CUPED,par jour}
\end{axis}
\end{tikzpicture}

\omcaption{Puissance pour détecter une économie de 0.3 point de base en fonction de la durée de l'expérience~: la formule à l'erreur type de chaque plan (traits) et la part
de 100 expériences simulées qui ont rejeté (points). La ligne pointillée vaut 80~\%. Données~: \texttt{rs\_abtest}.}\label{fig:rs:live-experiments:power}
\end{omfigure}

\paragraph{Consultations répétées.} Un desk qui consulte chaque matin le tableau de bord d'une expérience et l'arrête le premier jour où sa p-value passe sous 5~\% exécute un test différent
de celui que décrit la p-value. Sur 4\,000 expériences simulées sans aucun effet, un test à horizon fixe, consulté une seule fois, rejette dans 5~\% des cas~; consulté
chaque jour, il rejette au moins une fois dans 29~\% des cas sous 38 jours, dans 33~\% sous 60 jours et dans 38~\% sous 120 jours (\cref{fig:rs:live-experiments:peeking}).

Le \omterm{def:rs:live-experiments:mde}{test séquentiel} du rapport de vraisemblance par mélange (mixture sequential probability ratio test), issu du \omterm{def:rs:live-experiments:mde}{test séquentiel} de Wald (1945), rend la surveillance continue
valide (Johari, Koomen, Pekelis et Walsh, qui l'ont déployé sur une plateforme commerciale de tests). Pour une estimation $\hat\Delta$ de variance $V$ à la consultation courante, et des effets
tirés selon $N(0, \tau^2)$ sous l'hypothèse alternative, le rapport de vraisemblance par mélange contre l'absence d'effet vaut
\[
\Lambda = \sqrt{\frac{V}{V+\tau^2}}\,\exp\Bigl(\frac{\tau^2\hat\Delta^2}{2V(V+\tau^2)}\Bigr).
\]
Sous l'absence d'effet, $\Lambda$ est une martingale positive partant de 1, de sorte que la probabilité qu'elle dépasse à un instant quelconque $1/\alpha$ est au plus $\alpha$ (inégalité
de Ville). La p-value toujours valide est le minimum courant de $1/\Lambda$~; arrêter dès qu'elle passe sous $\alpha$ maintient l'erreur à $\alpha$, quel que soit le nombre
de consultations. Sur les mêmes expériences nulles, elle est passée sous 5~\% dans 46, 62 et 82 des 4\,000 expériences aux jours 38, 60 et 120. Avec l'économie réelle de 0.30 et \omterm{def:rs:live-experiments:cuped}{CUPED}, et
$\tau = 0.3$, la moitié des expériences s'arrêtent d'ici le jour 38, durée de l'horizon fixe, 4~\% d'ici le jour 10 et 98~\% d'ici le jour 120. Une modification nuisible de $+1$ point de base s'arrête
en médiane au bout de 6 jours. Le \omterm{def:rs:live-experiments:mde}{test séquentiel} renonce à une date de fin fixée à l'avance~; il conserve la taille du test sous surveillance et s'arrête tôt lorsque l'effet est grand, précisément quand
l'arrêt compte.

\begin{omfigure}
\begin{tikzpicture}
\begin{axis}[name=a,width=6.8cm,height=5.4cm,xlabel={\small jours de consultation quotidienne},ylabel={\small taux de faux positifs},tick label style={font=\footnotesize},
  xmin=0,xmax=120,ymin=0,ymax=0.42,grid=major,grid style={gray!25},title={\small absence d'effet},
  yticklabel style={/pgf/number format/fixed,/pgf/number format/precision=2},scaled y ticks=false,table/col sep=comma]
\addplot[black!60,thick] table[x=days,y=naive]{figdata/research/21-live-experiments/peeking.csv};
\addplot[omDef,very thick] table[x=days,y=msprt]{figdata/research/21-live-experiments/peeking.csv};
\draw[omThm,densely dotted,thick] (axis cs:0,0.05) -- (axis cs:120,0.05);
\end{axis}
\begin{axis}[at={(a.east)},anchor=west,xshift=1.3cm,width=6.8cm,height=5.4cm,xlabel={\small jour},ylabel={\small part arrêtée},
  tick label style={font=\footnotesize},xmin=0,xmax=120,ymin=0,ymax=1.02,grid=major,grid style={gray!25},title={\small économie de 0.30 pb},
  table/col sep=comma]
\addplot[omDef,very thick] table[x=days,y=share]{figdata/research/21-live-experiments/stopping.csv};
\draw[black!60,dashed] (axis cs:38,0) -- (axis cs:38,1.02);
\end{axis}
\end{tikzpicture}

\omcaption{À gauche~: la part des 4\,000 expériences nulles déclarées significatives à l'une quelconque des consultations quotidiennes, avec un test à horizon fixe (gris) et le \omterm{def:rs:live-experiments:mde}{test séquentiel}
du rapport de vraisemblance par mélange (bleu)~; la ligne pointillée vaut 5~\%. À droite~: le jour où le \omterm{def:rs:live-experiments:mde}{test séquentiel} s'arrête sous l'économie réelle, avec \omterm{def:rs:live-experiments:cuped}{CUPED}~; la ligne tiretée est l'horizon fixe de 38 jours
qui a une puissance de 80~\%. Données~: \texttt{rs\_abtest.peeking}, \texttt{stopping}.}\label{fig:rs:live-experiments:peeking}
\end{omfigure}

\paragraph{Le plan.} Tout ce qui détermine le résultat est écrit avant que le premier ordre soit assigné~: la métrique et ses garde-fous (taux d'exécution,
rejets, mark-outs), l'\omterm{def:rs:live-experiments:ab}{unité de randomisation}, la part, les covariables, l'horizon ou la règle séquentielle avec son $\tau$, l'\omterm{def:rs:live-experiments:mde}{effet minimal détectable},
et ce qui sera fait pour chaque issue. Le \omterm{def:rs:the-research-process:log}{journal de recherche} du chapitre~1 le consigne~; l'arithmétique des essais du chapitre~20 s'applique aux expériences tout autant qu'aux
\omterm{def:rs:vectorised-backtests:backtest}{backtests}, et un desk qui mène vingt expériences par an verra apparaître une fausse économie par an au seuil de 5~\%.

\section{Tutoriel~: le nouvel algorithme est-il meilleur~?}\label{tut:rs:live-experiments:tut}

\begin{tutorial}
\textbf{Objectif.} Simuler le flux d'ordres de la société avec un nouvel algorithme d'exécution, assigner les ordres par hachage, comparer les plans et les estimateurs, dimensionner l'expérience et
la surveiller de manière séquentielle. \textbf{État final~:} le tableau, les \cref{fig:rs:live-experiments:power,fig:rs:live-experiments:peeking}.
\begin{tutsteps}
\item \textbf{Assignation et estimation}~: le hachage, la différence de moyennes, \omterm{def:rs:live-experiments:cuped}{CUPED}.
\omcode{code/firm/abtest/firm_abtest.py}{35}{59}{Assignation par hachage, différence de moyennes et CUPED.}\label{lst:rs:live-experiments:assign}
\item \textbf{Le flux}~: \texttt{rs\_abtest.orders(days, seed, share, unit)} tire les ordres, leurs covariables, les apparentés et le résultat~; \texttt{hook()}
lance l'expérience de deux semaines à 10~\% et son contrôle SRM.
\item \textbf{Plans}~: \texttt{calibration()} lance 100 expériences par plan~; \texttt{rollout()} donne l'effet réel du passage à l'ensemble du flux.
\item \textbf{Surveillance séquentielle}~: le rapport de vraisemblance par mélange et la p-value toujours valide.
\omcode{code/firm/abtest/firm_abtest.py}{121}{129}{Le test séquentiel du rapport de vraisemblance par mélange.}\label{lst:rs:live-experiments:msprt}
\item \textbf{Lancer} \texttt{peeking()}, \texttt{stopping()} et \texttt{fig\_abtest.py}.
\end{tutsteps}
\textbf{Pour aller plus loin.} Faire en sorte que le nouvel algorithme nuise davantage à ses apparentés qu'il ne profite à lui-même (un débordement de 0.40) et observer un test au niveau de l'ordre approuver une
modification qui perd de l'argent~; randomiser par client plutôt que par action-jour et se demander quelle \omterm{def:rs:live-experiments:ab}{interférence} cela contient.
\end{tutorial}

\section{Construction~: la boîte à outils d'expérimentation}\label{bld:rs:live-experiments:abtest}

\begin{build}
\bfield{Rôle} Toute modification d'un composant de trading atteint le flux complet par un \omterm{def:rs:live-experiments:canary}{déploiement progressif}, et toute amélioration revendiquée provient d'une expérience randomisée,
analysée comme prévu.
\bfield{Interface} \texttt{bucket(unit, experiment, salt)}, \texttt{assign(units, experiment, share, salt)}, \texttt{diff\_means(y, t)},
\texttt{cuped(y, X, t)}, \texttt{stratified(y, t, strata)}, \texttt{cluster\_diff(y, t, cluster)}, \texttt{demean\_by(v, group)},
\texttt{power(effect, se, alpha)}, \texttt{mde(se, alpha, power)}, \texttt{sample\_size(effect, sd, alpha, power, share)},
\texttt{srm\_pvalue(n\_t, n\_c, share)}, \texttt{msprt(d, se, tau)}, \texttt{always\_valid\_p(d, se, tau)}.
\bfield{Règles} Assignation par hachage d'un identifiant stable, jamais par une table modifiable~; covariables fixées avant l'assignation~; l'\omterm{def:rs:live-experiments:ab}{unité de randomisation} contient
l'\omterm{def:rs:live-experiments:ab}{interférence}~; le contrôle SRM est effectué avant toute lecture d'une estimation~; un horizon fixe n'est pas consulté tôt, et une règle séquentielle est choisie avant le début.
\bfield{Tests d'acceptation} \texttt{code/firm/abtest/tests/}~: une assignation déterministe et équilibrée, indépendante d'une expérience à l'autre~; la variance de \omterm{def:rs:live-experiments:cuped}{CUPED} diminue de
$1-R^2$ sur une covariable connue, stratification entre les estimations brute et ajustée~; bras de grappes constants, ou une erreur~; puissance, \omterm{def:rs:live-experiments:mde}{effet minimal détectable} et
taille d'échantillon cohérents~; une p-value toujours valide non croissante et inférieure à 5~\% dans moins de 5~\% des 2\,000 flux nuls de 200 consultations, alors que les consultations quotidiennes
d'un test à horizon fixe dépassent 30~\%.
\bfield{Pour aller plus loin} Estimateurs pondérés par la variance pour des grappes inégales~; \omterm{def:rs:live-experiments:mde}{tests séquentiels} propres aux métriques de garde-fou~; plan en alternance avec effet de report.
\end{build}

\begin{omsources}
\item A.~Deng, Y.~Xu, R.~Kohavi et T.~Walker, ``Improving the sensitivity of online controlled experiments by utilizing pre-experiment data'', \emph{WSDM}, 2013.
\item R.~Kohavi, D.~Tang et Y.~Xu, \emph{Trustworthy Online Controlled Experiments}, Cambridge University Press, 2020.
\item A.~Fabijan et al., ``Diagnosing sample ratio mismatch in online controlled experiments'', \emph{KDD}, 2019.
\item R.~Johari, P.~Koomen, L.~Pekelis et D.~Walsh, ``Always valid inference: continuous monitoring of A/B tests'', \emph{Operations Research} 70(3), 2022
(et ``Peeking at A/B tests'', \emph{KDD}, 2017).
\item A.~Wald, ``Sequential tests of statistical hypotheses'', \emph{Annals of Mathematical Statistics} 16(2), 1945.
\item SEC, \emph{In the Matter of Knight Capital Americas LLC}, Release No.~34-70694, 16 octobre 2013.
\item I.~Bojinov, D.~Simchi-Levi et J.~Zhao, ``Design and analysis of switchback experiments'', \emph{Management Science} 69(7), 2023.
\item L.~W.~Miratrix, J.~S.~Sekhon et B.~Yu, ``Adjusting treatment effect estimates by post-stratification in randomized experiments'', \emph{Journal of the
Royal Statistical Society B} 75(2), 2013.
\end{omsources}

\section{Exercices}

\begin{exercise}[$\star$]\label{exo:rs:live-experiments:1}
L'expérience de deux semaines estime une économie de 0.55 point de base avec une erreur type de 0.48. Donner l'intervalle à 95~\% et le verdict.
\end{exercise}

\begin{exercise}[$\star$]\label{exo:rs:live-experiments:2}
La même expérience a prévu 10~\% de 24\,000 ordres et en a assigné 2\,485. Calculer la statistique SRM et sa p-value. Dire quelle décision prendre face à une p-value de 0.0001~?
\end{exercise}

\begin{exercise}[$\star$]\label{exo:rs:live-experiments:3}
L'erreur type d'une estimation est de 0.085 point de base. Calculer l'\omterm{def:rs:live-experiments:mde}{effet minimal détectable} à 5~\% et 80~\%, ainsi que la puissance contre une économie réelle de 0.3.
\end{exercise}

\begin{exercise}[$\star\star$]\label{exo:rs:live-experiments:4}
Indiquer lesquelles de ces grandeurs peuvent servir de covariables \omterm{def:rs:live-experiments:cuped}{CUPED} dans une expérience d'exécution~: l'estimation du coût avant transaction~; le mouvement du marché sur la fenêtre programmée de l'ordre~; la
durée réalisée de l'ordre~; le taux d'exécution de l'algorithme~; le shortfall de l'action la veille~? Avec des covariables expliquant 51~\% de la variance, calculer combien de
jours sur 77 restent~?
\end{exercise}

\begin{exercise}[$\star\star$]\label{exo:rs:live-experiments:5}
Un défaut coûte 5 points de base par ordre, l'écart type par ordre vaut 23.1, l'alarme est fixée à $z = 3$ et le flux est de 2\,400 ordres par jour. Combien
d'ordres, et combien de temps faut-il jusqu'à l'alarme, avec un canari à 1~\%, à 10~\% et à 50~\%~? Qu'est-ce qui aurait arrêté le défaut de Knight, et pourquoi cela ne relève pas des statistiques~?
\end{exercise}

\begin{exercise}[$\star\star$]\label{exo:rs:live-experiments:6}
Avec un gain direct de 0.30, un effet de débordement de 0.25 sur les ordres connexes entièrement traités et 88.3~\% des ordres ayant au moins un ordre connexe, quel est l'effet du déploiement complet~? Pourquoi la
randomisation par ordre mesure 0.30, et quelle unité mesure le déploiement~?
\end{exercise}

\begin{exercise}[$\star\star\star$]\label{exo:rs:live-experiments:7}
\emph{Programmation.} Lancer \texttt{rs\_abtest.stopping} avec l'économie réelle, pour $\tau = 0.1$, $\tau = 1$ ainsi que 0.3. Rapporter la médiane des jours d'arrêt et
expliquer pourquoi un petit et un grand $\tau$ ralentissent tous deux l'arrêt.
\end{exercise}

\begin{exercise}[$\star\star\star$]\label{exo:rs:live-experiments:8}
\emph{Trouver la faille.} «~Nous avons consulté le tableau de bord de l'expérience chaque matin, et le jour 23 la p-value est tombée à 0.04~; nous avons donc arrêté et déployé le nouvel
algorithme~: une économie significative à 5~\%.~»
\end{exercise}

\section{Problème~: le nouvel algorithme est-il meilleur~?}

\begin{problem}[{Problème du week-end --- une expérience d'exécution, conçue et lue}]\label{pb:rs:live-experiments:1}
Le flux simulé du chapitre~: 2\,400 ordres par jour, un nouvel algorithme qui économise 0.30 point de base sur ses propres ordres et coûte 0.25 à ses ordres connexes.

\medskip\noindent\textbf{Partie I --- La conception.}
\begin{enumerate}
\item Quelles sont la moyenne et l'écart type du shortfall d'un ordre, et quelles en sont les composantes~?
\item Quelles propriétés l'affectation par hachage procure-t-elle, et pourquoi comptent-elles sur un système de trading~?
\item Vérifier la répartition de l'expérience initiale.
\item Que mesure la randomisation par ordre, et combien le déploiement complet économise-t-il~?
\item Qu'est-ce que le \omterm{def:rs:live-experiments:canary}{déploiement progressif} de Knight aurait dû vérifier à chaque étape~?
\end{enumerate}

\medskip\noindent\textbf{Partie II --- Précision.}
\begin{enumerate}[resume]
\item Combien de jours, à la moitié du flux, la différence de moyennes exige-t-elle pour 0.3 point de base~?
\item Quelles covariables \omterm{def:rs:live-experiments:cuped}{CUPED} utilise-t-il, quel est leur $R^2$, et combien de jours restent-ils~?
\item Qu'apportent l'estimation avant transaction et la stratification qui s'appuie sur elle, et pourquoi si peu~?
\item Que coûte la randomisation par jour~?
\item Quelle est la précision du plan action-jour analysé à l'intérieur de la journée avec \omterm{def:rs:live-experiments:cuped}{CUPED}, et qu'estime-t-il~?
\end{enumerate}

\medskip\noindent\textbf{Partie III --- Arrêt.}
\begin{enumerate}[resume]
\item Quels taux de faux positifs donne le fait de consulter chaque jour un test à horizon fixe, aux jours 38 et 60~?
\item Combien des 4\,000 expériences nulles le \omterm{def:rs:live-experiments:mde}{test séquentiel} par mélange arrête-t-il aux mêmes jours~?
\item Avec l'économie réelle, quand le \omterm{def:rs:live-experiments:mde}{test séquentiel} s'arrête-t-il~?
\item En combien de temps détecte-t-il un changement néfaste de 1 point de base~?
\end{enumerate}

\medskip\noindent\textbf{Partie IV --- Le verdict.}
\begin{enumerate}[resume]
\item Énoncer le \emph{résultat nommé}~: le nombre de jours de trading nécessaires pour détecter une amélioration de 0.3 point de base à une puissance de 80~\%, avec et sans \omterm{def:rs:live-experiments:cuped}{CUPED}.
\item Le nouvel algorithme est-il meilleur~?
\item Comment le déployer~?
\item Combien de temps un canari à 1~\% mettrait-il à repérer un défaut coûtant 5 points de base par ordre, et sur quoi faudrait-il s'appuyer à la place~?
\item Que faut-il inscrire dans le plan de l'expérience avant le premier ordre~?
\item En une phrase~: que mesure un \omterm{def:rs:live-experiments:ab}{test A/B} sur le flux réel~?
\end{enumerate}
\end{problem}

\section{Questions d'entretien}

\begin{interviewq}[$\star$ \iqroles{trader, développeur}]\label{iq:rs:live-experiments:1}
Comment mettriez-vous en production, en toute sécurité, un nouvel algorithme d'exécution~?
\end{interviewq}

\begin{interviewq}[$\star\star$ \iqroles{chercheur}]\label{iq:rs:live-experiments:2}
Qu'est-ce que \omterm{def:rs:live-experiments:cuped}{CUPED}, quand aide-t-il, et qu'est-ce qui peut mal tourner avec lui~?
\end{interviewq}

\begin{interviewq}[$\star\star$ \iqroles{chercheur}]\label{iq:rs:live-experiments:3}
Pourquoi est-ce un problème de contrôler un \omterm{def:rs:live-experiments:ab}{test A/B} chaque jour et de l'arrêter dès qu'il est significatif, et que faites-vous à la place~?
\end{interviewq}

\begin{interviewq}[$\star\star$ \iqroles{chercheur, trader}]\label{iq:rs:live-experiments:4}
Un \omterm{def:rs:live-experiments:ab}{test A/B} randomisé par ordre a montré que le nouvel algorithme économisait 0.3 point de base. Après le déploiement complet, les coûts n'ont pas baissé. Que s'est-il passé~?
\end{interviewq}

\begin{interviewq}[$\star\star$ \iqroles{développeur}]\label{iq:rs:live-experiments:5}
Comment affectez-vous les ordres aux bras de l'expérience sur de nombreux serveurs, et comment vérifiez-vous que l'affectation a fonctionné~?
\end{interviewq}

\begin{interviewq}[$\star\star\star$ \iqroles{chercheur}]\label{iq:rs:live-experiments:6}
Dérivez la taille d'échantillon nécessaire pour détecter un effet $\Delta$ avec une puissance $1-\beta$. Avec un écart type de 23.1 points de base et 2\,400 ordres par jour, combien
de jours faut-il pour 0.3 point de base~?
\end{interviewq}
